\documentclass[aspectratio=169,10pt]{beamer}
\usetheme[numbering=fraction,progressbar=frametitle]{metropolis}
\usepackage{fontspec}
\setsansfont{Latin Modern Sans}
\usepackage{amsmath,amssymb,mathtools,graphicx,booktabs}
\usepackage{polyglossia}
\setdefaultlanguage{english}
\setbeamertemplate{navigation symbols}{}
\setbeamerfont{frametitle}{size=\large}
\setbeamerfont{normal text}{size=\normalsize}
\setlength{\parskip}{.45em}
\hypersetup{pdfauthor={},pdfsubject={Lecture slides},pdfkeywords={},pdftitle={Integration}}
\providecommand{\N}{\mathbb{N}}
\providecommand{\Z}{\mathbb{Z}}
\providecommand{\D}{\mathbb{D}}
\providecommand{\Q}{\mathbb{Q}}
\providecommand{\R}{\mathbb{R}}
\providecommand{\C}{\mathbb{C}}
\providecommand{\vect}[1]{\overrightarrow{#1}}
\providecommand{\vi}{\vec{\imath}}
\providecommand{\vj}{\vec{\jmath}}
\providecommand{\vk}{\vec{k}}
\providecommand{\norm}[1]{\left\lVert #1\right\rVert}
\providecommand{\vc}[2]{\begin{pmatrix}#1\\#2\end{pmatrix}}
\providecommand{\deter}[2]{\det\left(#1,#2\right)}
\providecommand{\abs}[1]{\left\lvert #1\right\rvert}
\providecommand{\dd}{\,\mathrm{d}}
\providecommand{\e}{\mathrm{e}}
\providecommand{\Card}{\operatorname{Card}}
\providecommand{\E}{\operatorname{E}}
\providecommand{\V}{\operatorname{V}}
\providecommand{\degre}{^{\circ}}
\providecommand{\un}[1]{\,\mathrm{#1}}
\providecommand{\sci}[2]{#1\times10^{#2}}
\providecommand{\pgcd}{\operatorname{PGCD}}
\providecommand{\ppcm}{\operatorname{PPCM}}
\title{Integration}
\subtitle{Lecture and worked examples}
\date{}
\author{}
\begin{document}
\begin{frame}[plain]
\vfill
{\usebeamerfont{title}\inserttitle\par}
\medskip
{\usebeamerfont{subtitle}\insertsubtitle\par}
\medskip
{\usebeamercolor[fg]{progress bar}\rule{\textwidth}{.5pt}}
\vfill
\end{frame}
\begin{frame}{Objectives}
Calculate and bound an integral; distinguish area from integral; use integration by parts; study a sequence of integrals.
\end{frame}
\begin{frame}{Area and antiderivative}
For \(f\) continuous and non-negative on \([a,b]\), \(a\le b\), the integral \(\int_a^b f(x)\,\dd x\) is the area under the curve, above the x-axis, between \(x=a\) and \(x=b\), expressed in area units.
\end{frame}
\begin{frame}{Area and antiderivative}
If \(f\) is continuous on an interval \(I\) and \(a\in I\), then \(F_a(x)=\int_a^x f(t)\,\dd t\) is the antiderivative of \(f\) that vanishes at \(a\). For any antiderivative \(F\),
\[\int_a^b f(x)\,\dd x=F(b)-F(a).\]
Every continuous function on an interval has antiderivatives. An integral is a number; an antiderivative is a function.
\end{frame}
\begin{frame}{Area and antiderivative}
For \(h>0\), the rectangles give \(hf(x)\le F_a(x+h)-F_a(x)\le hf(x+h)\). After dividing by \(h\) and taking the limit, the rate of change tends to \(f(x)\). For \(h<0\), the function being increasing gives \(f(x+h)\le f(t)\le f(x)\) on \([x+h,x]\). Integrating,
\[(-h)f(x+h)\le F_a(x)-F_a(x+h)\le(-h)f(x).\]
Dividing by \(-h>0\), we obtain
\[f(x+h)\le\dfrac{F_a(x+h)-F_a(x)}h\le f(x).\]
Both bounds tend to \(f(x)\). By continuity, \(F_a'=f\). Any other antiderivative is \(F_a+C\), hence the formula for evaluating between the limits.
\end{frame}
\begin{frame}{Properties}
For continuous functions and limits belonging to the same interval:
\[\int_a^a f=0,\quad\int_b^a f=-\int_a^b f,\quad
\int_a^c f=\int_a^b f+\int_b^c f,\quad
\int_a^b(\lambda f+\mu g)=\lambda\int_a^b f+\mu\int_a^b g.\]
If \(a\le b\) and \(f\le g\), then \(\int_a^b f\le\int_a^b g\). In particular, if \(m\le f\le M\), then \(m(b-a)\le\int_a^b f\le M(b-a)\).
For a function of arbitrary sign, the integral is defined using antiderivatives. The geometric area between the curve and the axis is \(\int_a^b|f|\), obtained by splitting at the sign changes. Between two curves, it equals \(\int_a^b|f-g|\).
\end{frame}
\begin{frame}{Properties}
\(\int_{-1}^{1}x\,\dd x=0\), but the area between the line and the axis over this interval is \(1\).
\end{frame}
\begin{frame}{Mean value and integration by parts}
For \(a<b\), the mean value is \(\mu=\dfrac1{b-a}\int_a^b f(x)\,\dd x\). The rectangle of height \(\mu\) and width \(b-a\) has the same signed area.
\end{frame}
\begin{frame}{Mean value and integration by parts}
If \(u,v\) are of class \(\mathcal C^1\) on \([a,b]\),
\[\int_a^b u(x)v'(x)\,\dd x=[u(x)v(x)]_a^b-\int_a^b u'(x)v(x)\,\dd x.\]
\end{frame}
\begin{frame}{Mean value and integration by parts}
Integrate \((uv)'=u'v+uv'\) then isolate the required integral. The bracket is indispensable.
\end{frame}
\begin{frame}{Mean value and integration by parts}
With \(u=x\) and \(v'=e^x\), \(\int_0^1xe^x\,\dd x=[xe^x]_0^1-\int_0^1e^x\,\dd x=e-(e-1)=1\).
\end{frame}
\begin{frame}{Rectangles and sequences of integrals}
For \(f\) continuous and increasing on \([a,b]\), with \(a<b\) and \(n\ge1\) and \(h=\dfrac{b-a}{n}\), the sums \(h\sum_{k=0}^{n-1}f(a+kh)\) and \(h\sum_{k=1}^{n}f(a+kh)\) bound the integral. Their difference is \(h(f(b)-f(a))\).
For \(I_n=\int_0^1x^nf(x)\,\dd x\), begin with the sign and comparisons of the integrand. If \(0\le f\le M\), then \(0\le I_n\le \dfrac{M}{n+1}\); the squeeze theorem gives \(I_n\to0\). Integration by parts can also provide a relation between \(I_n\) and \(I_{n+1}\).
\end{frame}
\begin{frame}{Rectangles and sequences of integrals}
\begin{columns}[T,onlytextwidth]
\begin{column}{.57\textwidth}
In \(I_n=\int_0^1x^nf(x)\,\dd x\), the expression \(x^0\) denotes the constant function \(1\), including at the endpoint \(0\). Thus \(I_0=\int_0^1f(x)\,\dd x\); no indeterminate calculation is needed.   The area between \(y=x\) and \(y=x^2\) on \([0,1]\) equals \(\dfrac{1}{6}\).
\end{column}
\begin{column}{.4\textwidth}
\centering
\includegraphics[width=\linewidth,height=4.5cm,keepaspectratio]{figures/terminale-11-1-en.pdf}
\end{column}
\end{columns}
\end{frame}
\begin{frame}{Research activity}
\textbf{Bounding an integral with rectangles - 35 min.}

We wish to bound \(I=\displaystyle\int_0^1x^2\,\dd x\). For an integer \(n\ge1\), we split \([0;1]\) into \(n\) intervals of equal length. Let \(L_n\) be the sum of the areas of the rectangles whose heights are the values of \(x^2\) at the left endpoints, and \(U_n\) the sum obtained at the right endpoints.
\end{frame}
\begin{frame}{Research activity (2)}
\begin{enumerate}
\setcounter{enumi}{0}\item For \(n=4\), draw both families of rectangles, then calculate \(L_4\) and \(U_4\).

\item Justify that \(L_4\le I\le U_4\).

\item Write a Python function with argument \(n\) that returns \(L_n\) and \(U_n\). Determine \(U_n-L_n\) algebraically.

\item Find the smallest integer \(n\) for which this bounding has a width at most equal to \(0.01\).

\item Calculate the exact value of \(I\) using an antiderivative and check that it lies within the bounds obtained for \(n=4\).

\end{enumerate}
\end{frame}
\end{document}
